\textbf{Tujuan:}\\
Formulir ini bertujuan untuk mengevaluasi kinerja dan kompetensi mahasiswa selama periode magang di perusahaan. Penilaian ini akan menjadi sarana pemberian kritik, saran, dan masukan yang bermanfaat bagi mahasiswa serta dapat digunakan sebagai dasar konversi SKS MBKM di universitas.

\vspace{0.3cm}
\noindent\textbf{Petunjuk Pengisian untuk Mentor:}\\
Bapak/Ibu Mentor yang terhormat,\\
Mohon kesediaannya untuk memberikan penilaian objektif terhadap kinerja mahasiswa magang yang tersebut di bawah ini. Penilaian dilakukan dengan memberikan nilai pada skala 1-5 pada setiap aspek serta menambahkan catatan atau komentar secara umum sebagai \textit{feedback}.

\vspace{0.5cm}
\subsection*{A. Data Mahasiswa dan Perusahaan}
\begin{table}[h]
	\renewcommand{\arraystretch}{1.5}
	\begin{tabular}{p{5.5cm} l p{8cm}}
		\textbf{Nama Mahasiswa}  & : & Andreandhiki Riyanta Putra                            \\
		\textbf{NIM}             & : & 23/517511/PA/22191                                    \\
		\textbf{Universitas}     & : & Universitas Gadjah Mada                               \\
		\textbf{Program Studi}   & : & S1 Ilmu Komputer                                      \\
		\textbf{Nama Perusahaan} & : & PT Vortex Buana Edumedia                              \\
		\textbf{Role/Posisi}     & : & Frontend Developer Intern                             \\
		\textbf{Nama Mentor}     & : & ..................................................... \\
		\textbf{Jabatan Mentor}  & : & ..................................................... \\
		\textbf{Periode Magang}  & : & 3 November 2025 - 3 Februari 2026                     \\
	\end{tabular}
\end{table}

\vspace{0.5cm}
\subsection*{B. Penilaian Aspek Kinerja}
\textit{(1 = Sangat Kurang, 2 = Kurang, 3 = Cukup, 4 = Baik, 5 = Sangat Baik)}

\vspace{0.3cm}
\renewcommand{\arraystretch}{1.5}
\begin{longtable}{|c|p{4.5cm}|c|p{6.5cm}|}
	\hline
	\textbf{No.} & \textbf{Aspek Penilaian}                       & \textbf{Nilai (1-5)} & \textbf{Kriteria / Deskripsi Penilaian}                                                                                                         \\
	\hline
	\endfirsthead

	\multicolumn{4}{c}%
	{{\bfseries Lanjutan dari halaman sebelumnya}}                                                                                                                                                                                         \\
	\hline
	\textbf{No.} & \textbf{Aspek Penilaian}                       & \textbf{Nilai (1-5)} & \textbf{Kriteria / Deskripsi Penilaian}                                                                                                         \\
	\hline
	\endhead

	\hline \multicolumn{4}{|r|}{{\textit{Berlanjut ke halaman berikutnya...}}}                                                                                                                                                             \\ \hline
	\endfoot

	\hline
	\endlastfoot

	\multicolumn{4}{|l|}{\textbf{MII21-3011 Internship: Pengembangan Fitur dan Modul Proyek}}                                                                                                                                              \\
	\hline
	1            & Implementasi \& Pengembangan Fitur             &                      & Mahasiswa mampu mengembangkan fitur dan modul proyek dengan baik sesuai dengan spesifikasi dan kebutuhan sistem.                                \\
	\hline
	2            & Kualitas Pelaporan \& Dokumentasi              &                      & Mahasiswa disiplin dalam mengisi logbook mingguan dan mampu menyusun Laporan Pengembangan Fitur dan Modul secara komprehensif.                  \\
	\hline
	3            & Presentasi Riset Inovasi                       &                      & Kemampuan mahasiswa dalam mempresentasikan hasil riset inovasi dan proses pengembangan modul yang telah dilakukan secara jelas dan terstruktur. \\
	\hline
	\multicolumn{4}{|l|}{\textbf{MII21-3015 Internship: Pengujian Unit dan Modul Proyek}}                                                                                                                                                  \\
	\hline
	4            & Pelaksanaan Pengujian (Testing)                &                      & Kemampuan mahasiswa dalam menyusun skenario dan melakukan pengujian unit serta modul proyek secara menyeluruh.                                  \\
	\hline
	5            & Laporan \& Evaluasi Pengujian                  &                      & Kualitas penyusunan laporan pengujian prototipe proyek serta konsistensi dokumentasi proses pengujian dalam logbook mingguan.                   \\
	\hline
	6            & Presentasi Magang \& Pengujian                 &                      & Kemampuan mempresentasikan hasil pengujian, temuan \textit{bug/error}, dan evaluasi magang secara komunikatif.                                  \\
	\hline
	\multicolumn{4}{|l|}{\textbf{MII21-4010 Soft Skill: Kemampuan Berkomunikasi}}                                                                                                                                                          \\
	\hline
	7            & Kecerdasan Emosional \& Empati                 &                      & Mahasiswa mampu memahami dan mengelola emosi untuk menghindari stres, mengatasi tantangan, serta menunjukkan empati kepada rekan kerja.         \\
	\hline
	8            & Komunikasi Efektif (Lisan \& Tulisan)          &                      & Mahasiswa mampu menyampaikan ide, progres pekerjaan, maupun hambatan dengan jelas dan ringkas.                                                  \\
	\hline
	9            & Sikap \& Keterbukaan                           &                      & Menunjukkan rasa percaya diri, rasa hormat terhadap sesama, serta memiliki pemikiran yang terbuka (\textit{open-minded}) dalam berdiskusi.      \\
	\hline
	10           & Mendengarkan Aktif (\textit{Active Listening}) &                      & Memiliki kemampuan mendengarkan instruksi maupun masukan dari mentor dan rekan kerja secara efektif.                                            \\
	\hline
	11           & Inisiatif \& Kualitas Pertanyaan               &                      & Mahasiswa proaktif dan mampu mengajukan pertanyaan yang baik, kritis, dan relevan terkait penyelesaian pekerjaan.                               \\
\end{longtable}

% \vspace{0.5cm}
\subsection*{C. Catatan Keseluruhan}
\textbf{Catatan Umum dan Saran Pengembangan:}\\
\textit{(Mohon berikan komentar terkait kelebihan mahasiswa, area yang perlu ditingkatkan, atau feedback konstruktif lainnya selama masa magang).}
\par\vspace*{0.2cm}
\noindent
\begin{tabular}{@{}p{\textwidth}@{}}
	\rule{\textwidth}{0.4pt} \\[0.1cm]
	\rule{\textwidth}{0.4pt} \\[0.1cm]
	\rule{\textwidth}{0.4pt} \\[0.1cm]
	\rule{\textwidth}{0.4pt}
\end{tabular}

\vspace{1cm}
\subsection*{D. Pengesahan}
Dengan ini saya menyatakan bahwa penilaian yang diberikan adalah benar dan sesuai dengan hasil pengamatan saya selama periode magang.

\vspace{1cm}
\begin{flushright}
	Yogyakarta, ........................... 2026 \\[2cm]
	\textbf{(.....................................................)} \\
	\textit{................}
\end{flushright}
